%% ch01 --- CPython, GIL i wersja free-threaded
%%
%% Napisane we wrzesniu 2026. Kazde twierdzenie w tym rozdziale pochodzi od
%% interpretera, a nie z pamieci: zobacz code/ch01/, ktorego trzy listingi CI
%% uruchamia przy kazdym pushu, oraz code/measure/e01_gil.py, czyli
%% eksperyment E1 -- zapisuje on surowy wynik do zacommitowanego pliku JSON,
%% bo czas zegarowy jest wlasnoscia maszyny i nie przejdzie bramki dryfu na
%% dwoch maszynach naraz.

\chapter{CPython, GIL i wersja free-threaded}
\label{ch:runtime}

\noindent
Masz już model środowiska uruchomieniowego. Proces hostuje jedno środowisko,
IL jest kompilowany z wyprzedzeniem, a potem jeszcze raz przez JIT przy
pierwszym wywołaniu, a pula wątków zamienia cztery rdzenie w czterokrotnie
więcej pracy. To dobry model, dzięki niemu jesteś produktywny \dash{} i trzy
z tych czterech zdań są tutaj nieprawdziwe. Nie przestają obowiązywać
jednocześnie: przestają pojedynczo, każde ma swoją nazwę i każde ma
jednolinijkowe pytanie, które możesz zadać interpreterowi stojącemu przed
tobą.

To jest metoda tego rozdziału i całej książki. Nic poniżej nie jest faktem,
który masz przyjąć na wiarę: każdy to linia, którą możesz uruchomić, a każdy
listing wypisuje swoją odpowiedź z działającego procesu, zamiast nieść ją
jako napis.

\section{Pięć pytań, na które interpreter odpowiada sam}
\label{sec:ch01-questions}
\index{CPython}\index{interpreter!introspekcja}

\noindent
Zacznij od pliku, który mówi, na czym stoisz.

\pyregion{code/ch01/interpreter.py}{facts}{Pięć pytań, na każde odpowiada działający proces}{lst:ch01-interpreter}

\noindent
Uruchom go, a przypięty interpreter odpowie:

\transcript{ch01-interpreter}

\noindent
Cztery z tych pięciu zmienią się w ciągu następnych czterech sekcji. Teraz
przeczytaj czwartą. \api{sys.\_jit} jest nowe w \pyver{}, mówi, że ten plik
wykonywalny został skompilowany \emph{z} kompilatorem just-in-time, i mówi,
że ten kompilator jest wyłączony. Obie połowy mają znaczenie. JIT istnieje,
więc kto twierdzi, że Python go nie ma, jest nieaktualny; jest domyślnie
wyłączony, więc kto rozumuje o twoich opóźnieniach na produkcji tak, jakby
JIT się rozgrzewał, myli się co do procesu, na który patrzy.

\begin{versionbox}
JIT w \pyver{} (PEP~744) jest eksperymentalny i domyślnie wyłączony, a noty
wydawnicze samego CPythona mówią, że nie jest zalecany na produkcji.
\code{PYTHON\_JIT=1} ustawia flagę włączenia; interpreter powyżej raportuje
wtedy \emph{dostępny} i \emph{włączony} razem. Plik wykonywalny można zbudować
tak, by niósł kompilator z wyłączoną flagą \dash{} i tym właśnie jest przypięty
interpreter, dlatego obie odpowiedzi się różnią. Ta książka nie zmierzyła, co
JIT robi z obciążeniem; uczciwie jest zacytować własne oszacowanie CPythona:
od nieco wolniej do nieco szybciej, zależnie od tego, co uruchamiasz.
\end{versionbox}

\begin{csbox}
\code{sys.\_is\_gil\_enabled()} i \code{sysconfig.get\_config\_var} to
najbliższy odpowiednik \code{RuntimeInformation.FrameworkDescription}: opisują
hosta, nie twój kod. Różnica polega na tym, że w Pythonie te odpowiedzi
niosą ciężar. Dwie z nich decydują o tym, czy pula wątków jest właściwym
narzędziem.
\end{csbox}

\section{Krok kompilacji istnieje i ma pamięć podręczną}
\label{sec:ch01-bytecode}
\index{bajtkod}\index{bajtkod!\_\_pycache\_\_}

\noindent
Zanim przeczytasz dalej, odpowiedz na to ze swojego obecnego modelu. Masz
\code{app.py}, który importuje \code{helper.py}. Uruchamiasz \code{python
app.py} dwa razy pod rząd. Ile razy kompilowany jest każdy z tych dwóch
plików?

\medskip
\noindent
Python kompiluje źródło do bajtkodu, zanim cokolwiek wykona, dokładnie tak
jak \code{csc} kompiluje do IL. Różnica polega na tym, kiedy to robi i gdzie
trzyma wynik. Kompilator uruchamia się przy imporcie, wynik ląduje obok
źródła w katalogu \code{\_\_pycache\_\_}, a \pkg{dis} z biblioteki
standardowej odczytuje go z powrotem.

\pyregion{code/ch01/bytecode.py}{function}{Cztery linie źródła, a kompilator ma coś do powiedzenia o każdej}{lst:ch01-function}

\noindent
\code{dis.dis} na tej funkcji wypisuje bajtkod, który wyprodukował
kompilator, z numerem linii źródła przy każdej instrukcji po lewej:

\transcript{ch01-bytecode}

\noindent
Odpowiedź na pytanie powyżej brzmi więc: \emph{dwa razy i raz}.
\code{helper.py} jest importowany, więc kompiluje się raz, a wynik trafia do
pamięci podręcznej; każde kolejne uruchomienie czyta ten cache.
\code{app.py} to plik, który podałeś w linii poleceń, a moduł podany w linii
poleceń nigdy nie jest cache'owany \dash{} kompiluje się ze źródła przy
każdym uruchomieniu. To niewielka praca, za którą płacisz za każdym razem, i
dlatego duży skrypt startuje odrobinę wolniej niż ten sam kod schowany za
importem.

\begin{trapbox}
\textbf{\enquote{Python nie ma kroku kompilacji, więc nie ma tu niczego jak
IL.}} Krok kompilacji istnieje, produkuje bajtkod, a bajtkod jest zapisany na
dysku, gdzie możesz go zobaczyć. Czego Python nie ma, to kroku kompilacji,
który musisz uruchomić, ani artefaktu budowania, który wysyłasz zamiast
źródła. Katalog \code{\_\_pycache\_\_} jest pamięcią podręczną i niczym
więcej: usuń go, a następne uruchomienie go odtworzy; wyślij go, a następny
interpreter może go zignorować; a \code{python -B} w ogóle go nie zapisuje.
\end{trapbox}

\mermaidfig{ch01-compile-step}{Te same dwa kroki w obu językach, w innym czasie i w inne miejsca}{od źródła do bajtkodu i do IL}

\begin{note}
Plik cache nosi nazwę interpretera, który go zapisał \dash{}
\valtext{e01.default.tag} w \pyver{} \dash{} więc dwie wersje mogą dzielić
katalog, nie czytając nawzajem swojego bajtkodu. Znacznik pochodzi z
\code{sys.implementation.cache\_tag}, a ścieżka z
\code{importlib.util.cache\_from\_source}, którą listing wypisuje, zamiast
składać ścieżkę samemu.
\end{note}

\begin{exercise}{e01_01_disassemble}{Przeczytaj bajtkod funkcji, którą napisałeś}
Zaimplementuj \code{instruction\_names}, które bierze funkcję i zwraca nazwy
opkodów, do jakich się skompilowała, w kolejności. Test deasembluje dwie
własne funkcje i sprawdza odpowiedź, więc musi ona pochodzić z obiektu
funkcji, a nie z listy, którą wpisałeś. Zwróć uwagę, czego test \emph{nie}
robi: nigdy nie wywołuje żadnej z tych funkcji. Bajtkod istnieje, odkąd moduł
został skompilowany.
\end{exercise}

\section{Moduł, asembly i po co jest \code{\_\_main\_\_}}
\label{sec:ch01-main}
\index{moduł}\index{sys.path@\texttt{sys.path}}

\noindent
Moduł to plik. Asembly to jednostka wdrożeniowa, która może nieść setki typów
z kilkudziesięciu plików, i niesie manifest. Prawie każda różnica w tej
sekcji bierze się z tej jednej.

Program można uruchomić na dwa sposoby i nie są one wymienne:

\begin{shellcmd}
python pkg/mod.py     # by path
python -m pkg.mod     # by module name
\end{shellcmd}

\noindent
Oba ustawiają \code{\_\_name\_\_} na \code{"\_\_main\_\_"} w module, który
uruchamiają. Różnią się tym, co trafia na ścieżkę importu. Pierwszy kładzie
tam katalog samego skryptu, więc na ścieżce jest \code{pkg/}, a
\code{import pkg} zawodzi. Drugi kładzie tam katalog bieżący, więc
\code{pkg} da się zaimportować, a moduł jest odnajdywany przez system
importu \dash{} i dlatego pod \code{-m} plik, który podałeś, dostaje wpis w
cache, a podany ścieżką nie dostaje.

\begin{csbox}
\code{python -m pkg.mod} to najbliższy odpowiednik \code{dotnet run}:
nazywasz rzecz, a nie jej plik, a środowisko ją rozwiązuje. \code{sys.path}
to sondowanie asemblatów, z tą różnicą, że jest zwykłą listą napisów, którą
możesz wypisać, zmienić i popsuć. Nie ma manifestu, nie ma silnej nazwy i nie
ma GAC; jest lista, a wygrywa pierwsze trafienie.
\end{csbox}

\begin{trapbox}
\textbf{\enquote{\code{if \_\_name\_\_ == "\_\_main\_\_"} to szablonowa
wstawka.}} To jedyna rzecz, która dzieli moduł dający się zaimportować od
modułu, który uruchamia swój program, gdy cokolwiek go zaimportuje. Zestaw
testów importuje twój plik, żeby dostać się do jednej funkcji; bez tej osłony
uruchamia też twoje \code{main}. Z osłoną ten sam plik jest i biblioteką, i
programem \dash{} dlatego każdy listing w tej książce ją ma.
\end{trapbox}

\begin{exercise}{e01_02_entrypoint}{Jeden plik, moduł i program}
Dokończ \code{main} i dodaj osłonę. Test importuje moduł i żąda, by nic się
nie wykonało, potem wywołuje \code{main} i żąda, by się wykonało, a na koniec
uruchamia ten sam plik jako program w podprocesie i żąda, by coś wypisał.
Jedna z tych trzech rzeczy zawodzi na pliku startowym z powodu, który warto
zauważyć: połowa z importem przechodzi na niedokończonym pliku, więc nie może
być testem samym w sobie.
\end{exercise}

\section{Jedna blokada i to, co naprawdę blokuje}
\label{sec:ch01-gil}
\index{GIL}\index{GIL!co szereguje}\index{wątki}

\noindent
Jeden proces CPythona to jeden interpreter z jedną blokadą. Wątek musi ją
trzymać, żeby wykonywać bajtkod, więc na domyślnej wersji dokładnie jeden
wątek wykonuje Pythona w danej chwili, niezależnie od tego, ile rdzeni
kupiłeś.

Zanim przyjdą liczby: cztery wątki, cztery rdzenie, każdy wątek liczy tę samą
czysto pythonową pętlę arytmetyczną. Ile, twoim zdaniem, zajmą te cztery
wątki w porównaniu z jednym wątkiem liczącym jedną pętlę?

\medskip
\noindent
Listing mierzy trzy rodzaje pracy na jednym wątku i na czterech. Oto te
obciążenia, a trzeciego ludzie nie przewidują:

\pyregion{code/ch01/workloads.py}{cpu}{Czysty Python: każdy krok tego trzyma blokadę}{lst:ch01-cpu}

\pyregion{code/ch01/workloads.py}{digest}{Ta sama praca co do kształtu, tyle że wykonana w rozszerzeniu C}{lst:ch01-digest}

\pyregion{code/ch01/workloads.py}{io}{Prawdziwy odczyt z gniazda, a nie uśpienie: wątek stoi w recv}{lst:ch01-io}

\noindent
Eksperyment E1 uruchamia ten plik na obu wersjach. Na wersji domyślnej, na
maszynie \valtext{e01.machine}, przy \val{e01.threads} wątkach:

\begin{center}
\small
\begin{tabularx}{\linewidth}{@{}Xrrr@{}}
\toprule
\textbf{Obciążenie} & \textbf{1 wątek} & \textbf{4 wątki} & \textbf{stosunek} \\
\midrule
\code{cpu}, czysty Python & \val{e01.default.cpu.one}\,s & \val{e01.default.cpu.many}\,s & \val{e01.default.cpu.ratio}$\times$ \\
\code{io}, odczyt z gniazda & \val{e01.default.io.one}\,s & \val{e01.default.io.many}\,s & \val{e01.default.io.ratio}$\times$ \\
\code{digest}, rozszerzenie C & \val{e01.default.digest.one}\,s & \val{e01.default.digest.many}\,s & \val{e01.default.digest.ratio}$\times$ \\
\bottomrule
\end{tabularx}
\end{center}

\noindent
Cztery wątki czystego Pythona zajmują \val{e01.default.cpu.ratio} razy tyle
co jeden. Nie pobiegły równolegle; pobiegły jeden po drugim i zapłaciły za
przełączanie. Pozostałe dwa wiersze nie kosztują nic, i to z tego samego
powodu w obu przypadkach: wątek stojący w \code{recv} nie wykonuje bajtkodu,
a wątek wewnątrz \pkg{hashlib} też nie \dash{} ta biblioteka zwalnia blokadę
na czas pracy i bierze ją z powrotem po niej.

\mermaidfig{ch01-one-lock}{Blokada jest trzymana, gdy biegnie bajtkod, i zwalniana wszędzie indziej}{co szereguje blokada}

\begin{trapbox}
\textbf{\enquote{Wątki w Pythonie są bezużyteczne, więc nie będę się nimi
zajmował.}} Dwa z trzech wierszy powyżej mówią co innego. Blokada szereguje
\emph{bajtkod}, a wątek, który czeka na gniazdo, plik albo bazę danych, lub
który jest w rozszerzeniu C i wykonuje pracę w C, nie trzyma niczego. Wątki
są właściwym narzędziem do pracy ograniczonej wejściem-wyjściem i
niewłaściwym do pracy ograniczonej procesorem \dash{} a to jest zdanie o
obciążeniu, nie o wątkach.
\end{trapbox}

\begin{note}
Konsekwencja tej blokady dla pętli zdarzeń \dash{} jedno blokujące wywołanie
w kodzie \code{async} zamraża wszystkie pozostałe korutyny, bez błędu i bez
wpisu w logu \dash{} należy do tomu towarzyszącego, \emph{LangChain,
LangGraph and Async Python}, który mierzy ją w swoim rozdziale~3. Ten
rozdział posiada model środowiska uruchomieniowego i na tym kończy.
\end{note}

\section{Wersja free-threaded to inny plik wykonywalny}
\label{sec:ch01-freethreaded}
\index{wersja free-threaded}

\noindent
Czytałeś, że \pyver{} usunął blokadę. Wróć do transkryptu w
\S\ref{sec:ch01-questions}: powstał na \pypatch{} i mówi, że blokada jest
włączona. Które z tych dwóch zdań jest nieprawdziwe?

\medskip
\noindent
Żadne. Free-threaded CPython to PEP~703, pojawił się jedną wersję wcześniej
jako eksperyment, a PEP~779 uczynił go wspieraną opcją budowania w
\pyver{}. \emph{Wspierana opcja budowania} to całe zdanie: są dwa pliki
wykonywalne. Ten, który instaluje prawie każdy i który ta książka przypina,
to wersja domyślna, z blokadą. Drugi to \code{python\pyver{}t}, budowany
osobno, instalowany osobno i rozpoznawalny po literze \code{t} w
\code{sys.abiflags} \dash{} stamtąd bierze się jego nazwa. Ten sam listing
uruchomiony pod każdym z nich daje:

\begin{center}
\small
\begin{tabularx}{\linewidth}{@{}Xll@{}}
\toprule
& \textbf{\code{python\pyver{}}} & \textbf{\code{python\pyver{}t}} \\
\midrule
raportowana wersja & \valtext{e01.default.version} & \valtext{e01.freethreaded.version} \\
wersja free-threaded & \valtext{e01.default.ft} & \valtext{e01.freethreaded.ft} \\
GIL włączony teraz & \valtext{e01.default.gil} & \valtext{e01.freethreaded.gil} \\
JIT dostępny / włączony & \valtext{e01.default.jit} & \valtext{e01.freethreaded.jit} \\
znacznik cache bajtkodu & \valtext{e01.default.tag} & \valtext{e01.freethreaded.tag} \\
\bottomrule
\end{tabularx}
\end{center}

\noindent
Przy dwóch wierszach warto się zatrzymać. Obie wersje raportują ten sam
znacznik cache, więc czytają i zapisują te same pliki
\code{\_\_pycache\_\_}: bajtkod jest identyczny, a różni się tylko środowisko
pod nim. A JIT nie jest na wersji free-threaded jedynie wyłączony \dash{} go
tam nie ma. To nie przypadek tego konkretnego pliku: noty wydawnicze
CPythona dla \pypatch{} mówią, że wersje free-threaded nie wspierają
kompilacji JIT, więc te dwa eksperymenty \pyver{} nie składają się ze sobą.

Teraz te same trzy obciążenia, na \code{python\pyver{}t}:

\begin{center}
\small
\begin{tabularx}{\linewidth}{@{}Xrrr@{}}
\toprule
\textbf{Obciążenie} & \textbf{1 wątek} & \textbf{4 wątki} & \textbf{stosunek} \\
\midrule
\code{cpu}, czysty Python & \val{e01.freethreaded.cpu.one}\,s & \val{e01.freethreaded.cpu.many}\,s & \val{e01.freethreaded.cpu.ratio}$\times$ \\
\code{io}, odczyt z gniazda & \val{e01.freethreaded.io.one}\,s & \val{e01.freethreaded.io.many}\,s & \val{e01.freethreaded.io.ratio}$\times$ \\
\code{digest}, rozszerzenie C & \val{e01.freethreaded.digest.one}\,s & \val{e01.freethreaded.digest.many}\,s & \val{e01.freethreaded.digest.ratio}$\times$ \\
\bottomrule
\end{tabularx}
\end{center}

\noindent
Pierwszy wiersz to cały eksperyment. Cztery wątki czystego Pythona kosztują
teraz \val{e01.freethreaded.cpu.ratio} raza tyle co jeden wątek, zamiast
\val{e01.default.cpu.ratio}, a ta sama praca kończy się
\val{e01.cpu.wallclock.gain} raza szybciej w czasie zegarowym. W kodzie nie
zmieniło się nic; zmienił się interpreter.

\begin{warning}
Czytaj obie tabele z rozrzutem w pamięci. Żaden stosunek nie trzymał się tu
stabilniej niż mniej więcej \val{e01.spread} procent w
\val{e01.repeats} przebiegach stojących za każdym wierszem, więc w tych
sześciu wierszach są dwie odpowiedzi, a nie sześć: wszystko w okolicach
\val{e01.freethreaded.cpu.ratio} i jeden wiersz z
\val{e01.default.cpu.ratio}. Czasy bezwzględne należą do maszyny
\valtext{e01.machine} w dniu \valtext{e01.recorded}; to, co się przenosi, to
stosunki, bo wynikają z blokady, a nie ze sprzętu.

Jest tu jedna liczba, którą ten eksperyment podaje i z której odmawia
wyciągania wniosków. Ile free-threading kosztuje pojedynczy wątek? O to każdy
pyta najpierw. Obie kolumny jednowątkowe różnią się o
\val{e01.cpu.singlethread.diff} procent, a noty wydawnicze samego CPythona
dla \pypatch{} szacują karę jednowątkową na mniej więcej pięć do dziesięciu
procent, zależnie od platformy i kompilatora \dash{} więc jedno zgadza się z
drugim, a ta zgodność nic nie znaczy, bo \val{e01.cpu.singlethread.diff}
procent przy rozrzucie \val{e01.spread} procent nie jest pomiarem. W kolejnych
zapisach na tej maszynie ta sama różnica wychodziła jeden, siedem i jedenaście
procent. Spokojniejsza maszyna i obciążenie zaprojektowane wyłącznie pod to
pytanie by to rozstrzygnęły; póki nikt tego nie uruchomił, cytowana liczba
jest ich, nie tej książki.
\end{warning}

\mermaidfig{ch01-two-binaries}{Jeden numer wersji, dwa pliki wykonywalne i jeden cache bajtkodu między nimi}{dwie wersje, jeden cache bajtkodu}

\begin{trapbox}
\textbf{\enquote{\pyver{} jest teraz free-threaded, więc GIL zniknął.}}
Zniknął w wersji, o którą trzeba poprosić. \code{uv python install
\pyver{}t} ją pobiera, \code{python\pyver{}t} ją uruchamia, a dopóki nie
wpiszesz jednego z tych poleceń, każdy proces na maszynie wciąż ma blokadę.
Wersja i proces to też osobne pytania: uruchom \code{python\pyver{}t} z
\code{PYTHON\_GIL=1}, a dostaniesz plik free-threaded z blokadą włączoną z
powrotem \dash{} dlatego listing w \S\ref{sec:ch01-questions} pyta o jedno i
o drugie.
\end{trapbox}

\section{Decyzja, którą to kupuje}
\label{sec:ch01-decision}
\index{wątki!wybór puli}

\noindent
Cztery wiersze pomiaru zwijają się w jedną regułę, którą stosujesz bez
mierzenia czegokolwiek ponownie: \textbf{blokada kosztuje cię dokładnie
wtedy, gdy kilka wątków naraz wykonuje bajtkod Pythona}. Wszystko inne
\dash{} czekanie na sieć, czekanie na dysk, czekanie na bazę danych albo
spędzanie czasu w bibliotece napisanej w C \dash{} jest już równoległe i było
takie przed \pyver{}.

Jest więc jedna komórka tabeli, w której wątki są złą odpowiedzią, a na
wersji domyślnej odpowiada się na nią procesami, z których każdy ma własną
blokadę. To właśnie ostatnie ćwiczenie każe ci zapisać.

\begin{exercise}{e01_03_which_pool}{Wątki czy procesy, i na której wersji}
Zaimplementuj \code{pick\_pool}, które bierze jeden z trzech rodzajów
obciążenia i zwraca \code{"threads"} albo \code{"processes"}. Wersję bierze
jako parametr, domyślnie taką, jaką raportuje działający interpreter, więc
test może zadać ci obie, choć masz zainstalowaną jedną. Tylko jedna z sześciu
komórek mówi \enquote{procesy}, a ćwiczenie polega na tym, żeby wiedzieć
która.
\end{exercise}

\noindent
Dwa wątki tej reguły podjęte są później. Rozdział~\ref{ch:asyncio} wykonuje
tę samą pracę jednym wątkiem i pętlą zdarzeń, co jest inną odpowiedzią na to
samo pytanie o czekanie. A powód, dla którego proces jest kosztownym wyjściem
awaryjnym \dash{} że to osobny interpreter z osobną pamięcią i kosztem
serializacji na granicy \dash{} jest pytaniem o pakowanie, zanim stanie się
pytaniem o współbieżność, i od tego zaczyna się następny rozdział.
